\par\Needspace{6\baselineskip}\medskip
\hypertarget{native-149-article108}{}\pdfbookmark[2]{108.}{article-108}
\begin{center}108.\end{center}
If the parabola is regarded as an ellipse whose major axis is infinitely large, the expression for the time found in the preceding article becomes $\dfrac1{6k}\{(r+r'+\rho)^{\frac32}\mp(r+r'-\rho)^{\frac32}\}$. Since this derivation of the formula may perhaps seem open to some doubts, we give another derivation that is independent of the ellipse.
\par\SourceLeaf{150}
Setting for brevity
\[
\operatorname{tang}\tfrac12v=\theta,\qquad\operatorname{tang}\tfrac12v'=\theta',
\]
\noindent we obtain
\[
\begin{aligned}
r&=\tfrac12p(1+\theta\theta),&r'&=\tfrac12p(1+\theta'\theta'),\\[.6em]
\cos v&=\frac{1-\theta\theta}{1+\theta\theta},&\cos v'&=\frac{1-\theta'\theta'}{1+\theta'\theta'},\\[.6em]
\sin v&=\frac{2\theta}{1+\theta\theta},&\sin v'&=\frac{2\theta'}{1+\theta'\theta'}\mathpunct{\cdot}
\end{aligned}
\]
\noindent This gives
\[
\begin{aligned}
r'\cos v'-r\cos v&=\tfrac12p(\theta\theta-\theta'\theta')\\[.6em]
r'\sin v'-r\sin v&=p(\theta'-\theta),
\end{aligned}
\]
\noindent and hence
\[
\rho\rho=\tfrac14pp(\theta'-\theta)^2(4+(\theta'+\theta)^2).
\]
\noindent It is readily seen that
\[
\theta'-\theta=\frac{\sin f}{\cos\tfrac12v\cos\tfrac12v'},
\]
\noindent is a positive quantity. Setting
\[
\BareRoot{(1+\tfrac14(\theta'+\theta)^2)}=\eta,
\]
\noindent then gives
\[
\rho=p(\theta'-\theta)\eta.
\]
\noindent We also have
\[
r+r'=\tfrac12p(2+\theta\theta+\theta'\theta')=p(\eta\eta+\tfrac14(\theta'-\theta)^2):
\]
\noindent and therefore
\[
\begin{aligned}
\frac{r+r'+\rho}{p}&=(\eta+\tfrac12(\theta'-\theta))^2\\[.6em]
\frac{r+r'-\rho}{p}&=(\eta-\tfrac12(\theta'-\theta))^2.
\end{aligned}
\]
\HistoricalUnmarkedNote{[Manuscript note:] $\eta$ is the cosecant of the angle between the chord and the axis\\
$\tfrac12(\theta'+\theta)$ is its cotangent\\
$\dfrac{r'-r}{\rho}$ is its cosine $=\cos\psi$
\[
\begin{aligned}
p&=\tfrac12\sin\psi^2\{r+r'+\BareRoot{((r+r')^2-\rho\rho)}\}\\[.6em]
&=(r'+r)\frac{\rho\rho-(r'-r)^2}{\rho\rho}\Bigl\{1-\tfrac14\frac{\rho\rho}{(r'+r)^2}\\[.6em]
&\qquad-\tfrac1{16}\frac{\rho^4}{(r'+r)^4}-\cdots\Bigr\}.
\end{aligned}
\]}
\par
The first equation gives at once
\[
+\BareRoot{\frac{r+r'+\rho}{p}}=\eta+\tfrac12(\theta'-\theta),
\]
\noindent since $\eta$ and $\theta'-\theta$ are positive quantities. Now $\tfrac12(\theta'-\theta)$ is less \SourceLeaf{151}or greater than $\eta$, according as
\[
\eta\eta-\tfrac14(\theta'-\theta)^2=1+\theta\theta'=\frac{\cos f}{\cos\tfrac12v\cos\tfrac12v'}
\]
\noindent is positive or negative. From the second equation we must therefore conclude that
\[
\pm\BareRoot{\frac{r+r'-\rho}{p}}=\eta-\tfrac12(\theta'-\theta),
\]
\noindent where the upper or lower sign is to be taken according as the angle swept out about the Sun is below $180^0$ or above $180^0$.
\par
From the equation following the second in Article 98, we also have
\[
\begin{aligned}
\frac{2kt}{p^{\frac32}}&=(\theta'-\theta)(1+\theta\theta'+\tfrac13(\theta'-\theta)^2)\\[.6em]
&=(\theta'-\theta)(\eta\eta+\tfrac1{12}(\theta'-\theta)^2)\\[.6em]
&=\tfrac13(\eta+\tfrac12(\theta'-\theta))^3-\tfrac13(\eta-\tfrac12(\theta'-\theta))^3,
\end{aligned}
\]
\begingroup\postdisplaypenalty=10000
\noindent from which it follows directly that
\[
kt=\tfrac16\{(r+r'+\rho)^{\frac32}\mp(r+r'-\rho)^{\frac32}\}
\]
\noindent with the upper or lower sign according as $2f$ is below or above $180^0$.\endgroup
\HistoricalUnmarkedNote{[Manuscript note:] The mean time between two positions on a parabola is given by the formula
\[
\begin{aligned}
&\frac1{2k\BareRoot{18}}\Biggl\{-\Biggl(\frac{r'-r}{\BareRoot{\frac{r'+r+\rho}{2}}\mp\BareRoot{\frac{r'+r-\rho}{2}}}\Biggr)^3\\[.8em]
&\qquad+\frac{3(r'r'-rr)}{\BareRoot{\frac{r'+r+\rho}{2}}\mp\BareRoot{\frac{r'+r-\rho}{2}}}\Biggr\}\\[.8em]
&=-\frac1{3k}\Biggl(\frac{r'-r}{\BareRoot{(r'+r+\rho)}\mp\BareRoot{(r'+r-\rho)}}\Biggr)^3\\[.8em]
&\qquad+\frac1{2k}\frac{r'r'-rr}{\BareRoot{(r'+r+\rho)}\mp\BareRoot{(r'+r-\rho)}}\mathpunct{\cdot}
\end{aligned}
\]
\par\noindent Perihelion distance $=\dfrac{\rho\rho-(r'-r)^2}{4\{r'+r\mp\BareRoot{((r'+r)^2-\rho\rho)}\}}\mathpunct{\cdot}$}
